\chapter{Groups}
\section{Groups}

\begin{example}
  The \emph{general linear group} \[
    GL_n(\R) = \{A\in M_n(\R) \mid \det(A)\neq 0\}
  \] is a group under matrix multiplication.
\end{example}

\begin{example}
  Let \(X\) be a nonempty set, then the \emph{symmetric group} \[
    S_X = \{\sgm:X\to X \mid \sgm\text{ is a bijection}\}
  \] is a group under composition of functions. In particular, we put \(S_n = S_{\{1,2,\ldots,n\}}\).
\end{example}
\section{Subgroups}

\begin{lemma}[subgroup criterion]
  A subset \(H\) of a group \(G\) is a subgroup if and only if \begin{enumerate}
    \item \(H\) is nonempty,
    \item \(ab^{-1}\in H\) for all \(a,b\in H\).
  \end{enumerate}
\end{lemma}

\begin{example}
  The \emph{special linear group} \[
    SL_n(\R) = \{A\in GL_n(\R) \mid \det(A) = 1\}
  \] is a subgroup of \(GL_n(\R)\). The \emph{orthogonal group} \[
    O_n(\R) = \{A\in GL_n(\R) \mid A^TA = I_n\}
  \] is a subgroup of \(GL_n(\R)\). The \emph{special orthogonal group} \[
    SO_n(\R) = \{A\in O_n(\R) \mid \det(A) = 1\}
  \] is a subgroup of \(O_n(\R)\).
\end{example}

\begin{example}
  The \emph{unitary group} \[
    U(n) = \{A\in GL_n(\C) \mid A^*A = I_n\}, \quad
  \] is a subgroup of \(GL_n(\C)\). The \emph{special unitary group} \[
    SU(n) = \{A\in U(n) \mid \det(A) = 1\}
  \] is a subgroup of \(U(n)\).

  In particular, we have \[
    U(1) = \{z\in\C \mid |z| = 1\}.
  \]
\end{example}

\begin{example}
  The \emph{alternating group} \[
    A_n = \{\sgm\in S_n \mid \sgm\text{ is even}\}
  \] is a subgroup of \(S_n\).
\end{example}

\begin{proposition}
  Consider the embedding \[
    \iota:GL_n(\R)\longrightarrow GL_{n+1}(\R),
    \quad A\longmapsto \begin{pmatrix}
      A & 0 \\
      0 & 1
    \end{pmatrix},
  \] then \(GL_n(\R)\cong \iota(GL_n(\R)) < GL_{n+1}(\R)\). 
\end{proposition}

\begin{proposition}
  Let \((H_\a)_\a\) be a family of subgroups of \(G\). Then \(\bigcap_{\a\in I}H_\a\) is a subgroup of \(G\).
\end{proposition}
\begin{proof}
  We check the subgroup criterion.
  \begin{itemize}
    \item \(e\in\bigcap_{\a\in I}H_\a\), since \(e\in H_\a\) for all \(\a\in I\).
    \item Pick \(a,b\in\bigcap_{\a\in I}H_\a\), then \(ab^{-1}\in H_\a\) for all \(\a\in I\), so \(ab^{-1}\in\bigcap_{\a\in I}H_\a\).\qedhere
  \end{itemize}
\end{proof}

\begin{proposition}
  Let \(H_1\leq H_2\leq H_3\leq\cdots\) be an ascending chain of subgroups of a group \(G\). Then \(\bigcup_{n=1}^\infty H_n\) is a subgroup of \(G\). 
\end{proposition}
\begin{proof}
  We check the subgroup criterion.
  \begin{itemize}
    \item \(e\in\bigcup_{n=1}^\infty H_n\), since \(e\in H_1\).
    \item Pick \(a,b\in\bigcup_{n=1}^\infty H_n\), then we have \(a\in H_i\) and \(b\in H_j\) for some \(i,j\in\N\), so \(a,b\in H_{\max\{i,j\}}\).
    Now we have \[
      ab^{-1}\in H_{\max\{i,j\}}\subset \bigcup_{n=1}^\infty H_n. \qedhere
    \]
  \end{itemize}
\end{proof}


\section{Cyclic Groups}

\section{Dihedral Groups}

Symmetries of a regular \(n\)-gon on a plane.

\begin{proposition}
  For \(n\geq 3\), \[
    D_n = \<r,s \mid r^n, s^2, (rs)^2\>.
  \]
\end{proposition}

\begin{proposition}
  In \(D_n\), we have \(sr  = r^{-1}s\). Consequently, \(sr^k = r^{-k}s\).
\end{proposition}

\section{Symmetric Groups}

Rigid motions of a regular \(n\)-simplex in \(\R^{n-1}\).

\newpage

\begin{proposition}\label{prop:S_n_transp_conjg}
  Let \(\sgm\in S_n\) and \(\tau = (a_0\,a_1\,\ldots\,a_{k-1})\) be a \(k\)-cycle, then \[
    \sgm(a_0\,a_1\,\ldots\,a_{k-1})\sgm^{-1} = (\sgm(a_0)\ \sgm(a_1)\,\ldots\,\sgm(a_{k-1})).
  \]
\end{proposition}
\begin{proof}
  Let \(\sgm(a_i)\), where \(i\in\Z_k\) then
  \[
    \sgm(a_i) \xmapsto[\hspace*{2em}]{\sgm^{-1}} a_i \xmapsto[\hspace*{2em}]{\tau} a_{i+1} \xmapsto[\hspace*{2em}]{\sgm} \sgm(a_{i+1}). \qedhere
  \]
\end{proof}

\begin{example}
  \(\Z_n \cong \big\<(1\ 2\cdots\ n)\big\> < S_n\).
\end{example}

\begin{example}
  \(D_n \cong \<r,s\> < S_n\) where \[
    r = \begin{pmatrix}
      1 & 2 & \cdots & n \\
      2 & 3 & \cdots & 1
    \end{pmatrix}, \quad
    s = \begin{pmatrix}
      1 & 2 & \cdots & n \\
      n & n-1 & \cdots & 1
    \end{pmatrix}.
  \]
\end{example}

Recall that \[
  V_4 = \Z_2\times\Z_2 \cong \<x,y \mid x^2 = y^2 = e, xy = yx\>.
\]
\begin{example}
  \(V_4 \cong \big\<(1\ 2),(3\ 4)\big\> = \{e,(1\ 2),(3\ 4),(1\ 2)(3\ 4)\}\).
\end{example}
\begin{proof}
  We can realize \(V_4\) as a subgroup of \(S_4\) by \[
    x \mapsto (1\ 2), \quad
    y \mapsto (3\ 4). \qedhere
  \]
\end{proof}

\begin{example}
  \(V_4 \cong \big\<(1\ 2)(3\ 4),(1\ 3)(2\ 4)\big\> = \{e,(1\ 2)(3\ 4),(1\ 3)(2\ 4),(1\ 4)(2\ 3)\}\).
\end{example}
\begin{proof}
  We can realize \(V_4\) as a subgroup of \(S_4\) by \[
    x \mapsto (1\ 2)(3\ 4), \quad
    y \mapsto (1\ 3)(2\ 4). \qedhere
  \]
\end{proof}

\begin{example}
  \(\Z_m\times\Z_n \cong \big\<(1\ 2\cdots\ m),(m+1\ m+2\cdots\ m+n)\big\> < S_{m+n}\).
\end{example}

\section{Quaternions}